\section{Ice Lattice Energies and Relative Phase Energies}
\label{sec:ice-basis}

All thirteen phases from DMC-ICE13 have completed AE calculations with both orbital bases: Ih, II, III, IV, VI, VII, VIII, IX, XI, XIII, XIV, XV, and XVII. Their electronic lattice energies provide a second test of the basis sensitivity observed for the molecular crystals. The DMC references are from Della Pia \emph{et al.}\cite{DellaPia2022Ice} Every phase is included at both basis levels; no phase-dependent basis selection is made. The last QZVPP calculation, XIII, converged in 41 SCF steps with a residual of $3.8\times10^{-7}$ and a normal program exit. Its TZVPP counterpart converged in 37 steps with a residual of $4.3\times10^{-7}$. Both exact execution records are archived.

Both basis levels use the same fixed crystal geometries, Skala--D3(BJ), explicit-electron Hamiltonian, joint one-centre reconstruction, 800/60~Ry cutoffs, and 150/770 atom quadrature. The SCF threshold is $5\times10^{-7}$ with OT. The QZVPP crystals use Gamma-centred $3^3$ meshes with full point-group reduction through spglib. All executions use one process and 32 CPU threads on ROSI. The executed inputs specify the symmetry tolerances and irreducible meshes for each phase.

The earlier TZVPP calculations use CP2K's K290 symmetry backend, including its inversion fallback for ice IX. Separate same-basis spglib checks for IX and VII change the energy by $-0.000035$ and $+0.000200$~kJ~mol$^{-1}$ per molecule, respectively. These controls distinguish the symmetry-backend change from the much larger basis shifts. For XIII, the irreducible meshes contain 14 points at TZVPP and 10 at QZVPP; neither execution uses the inversion fallback. Its inputs differ only in the orbital basis, project name, and explicit symmetry-backend/tolerance settings. The two same-basis checks are representative controls, not an independent backend test for every phase.

The gas-phase reference is a centred isolated water molecule at the reconstructed Partridge--Schwenke equilibrium geometry, with OH distance 0.95865~\AA{} and HOH angle $104.348^\circ$.\cite{Gillan2012WaterClusters} Its total energies are $-76.412709030115$ and $-76.436220026726$~hartree at TZVPP and QZVPP, respectively. This geometrical convention follows the equilibrium molecular reference, but the original DMC monomer coordinate file was not available for a direct file-level comparison. Any molecular-reference mismatch would shift all absolute ice lattice energies together. Relative phase energies cancel that molecular reference exactly.

\begin{table}[htbp]
\centering\small
\caption{Preliminary six-phase ice basis comparison. Electronic lattice energies and signed QZVPP--DMC residuals are in kJ~mol$^{-1}$ per water molecule. Parentheses give DMC statistical uncertainties. The averages concern these same six phases only, not the complete DMC-ICE13 benchmark.}
\label{tab:ice-basis}
\begin{tabular}{lrrrr}
\toprule
Phase & TZVPP & QZVPP & DMC & QZVPP--DMC\\
\midrule
VI & -80.663 & -56.275 & -57.67(7) & +1.395\\
VII & -80.495 & -55.128 & -54.46(7) & -0.668\\
VIII & -79.642 & -54.803 & -55.22(8) & +0.417\\
IX & -78.841 & -57.209 & -58.85(7) & +1.641\\
XV & -80.474 & -56.286 & -57.71(7) & +1.424\\
XVII & -75.684 & -56.674 & -57.70(8) & +1.026\\
\midrule
MAE & 22.365 & 1.095 & -- & --\\
\bottomrule
\end{tabular}
\end{table}


On all thirteen phases, the absolute lattice-energy MAE decreases from 21.349 to 1.161~kJ~mol$^{-1}$ (Table~\ref{tab:ice-basis}), a reduction of approximately 95\%. With the signed residual defined as calculated minus DMC energy, the mean signed error changes from $-21.349$ to $+1.058$~kJ~mol$^{-1}$, and the RMSE decreases from 21.511 to 1.292~kJ~mol$^{-1}$. All thirteen TZVPP energies overbind. Twelve QZVPP energies are less binding than DMC, while VII is more binding by 0.668~kJ~mol$^{-1}$; the largest absolute QZVPP residual is $+2.332$~kJ~mol$^{-1}$ for III. Thus the large systematic TZVPP overbinding is not retained with the enlarged AE basis and fully reduced k-point protocol. No empirical energy shift is applied. As for the three molecular crystals, the enlarged AE orbital space is essential to the balance between crystal and gas-phase energies. This result does not establish a complete-basis or counterpoise-corrected limit. The remaining errors generally exceed the quoted DMC statistical uncertainties of 0.07--0.08~kJ~mol$^{-1}$; those error bars are not a complete estimate of DMC systematic uncertainty.

For XIII specifically, the TZVPP and QZVPP lattice energies are $-79.942$ and $-57.036$~kJ~mol$^{-1}$, compared with the DMC value of $-57.33(7)$~kJ~mol$^{-1}$. The signed error therefore changes from $-22.612$ to $+0.294$~kJ~mol$^{-1}$. Its QZVPP absolute agreement is among the best in the set, but does not by itself establish an accurate relative stability.

\begin{table}[htbp]
\centering\small
\caption{Relative electronic energies with respect to ice Ih, in kJ~mol$^{-1}$ per water molecule. The MAE excludes the identically zero Ih row. DMC differences are calculated from the rounded absolute energies in Table II of Ref.~\citenum{DellaPia2022Ice}, so the XI value is 0.16 rather than the separately rounded 0.15 in its Table I. The gas-phase reference cancels exactly.}
\label{tab:ice-relative}
\begin{tabular}{lrrrr}
\toprule
Phase & TZVPP & QZVPP & DMC & QZVPP--DMC\\
\midrule
Ih & 0.000 & 0.000 & 0.00 & 0.000\\
II & -3.881 & -0.190 & 0.31 & -0.500\\
III & 0.133 & 1.859 & 1.25 & 0.609\\
IV & -0.592 & 2.918 & 3.83 & -0.912\\
VI & -3.716 & 1.452 & 1.78 & -0.328\\
VII & -3.547 & 2.600 & 4.99 & -2.390\\
VIII & -2.694 & 2.924 & 4.23 & -1.306\\
IX & -1.894 & 0.518 & 0.60 & -0.082\\
XI & -0.095 & -0.033 & 0.16 & -0.193\\
XIV & -4.068 & 0.585 & 1.70 & -1.115\\
XV & -3.527 & 1.442 & 1.74 & -0.298\\
XVII & 1.263 & 1.054 & 1.75 & -0.696\\
\midrule
MAE & 4.087 & 0.766 & -- & --\\
\bottomrule
\end{tabular}
\end{table}



The relative energies in Table~\ref{tab:ice-relative} test the balance among condensed phases without a gas-phase reference. For a phase $p$ containing $N_p$ water molecules per cell, we evaluate
\begin{equation}
 \Delta E_p = E_p/N_p-E_{\mathrm{Ih}}/N_{\mathrm{Ih}}.
\end{equation}
Across all twelve nontrivial differences, the MAE decreases from 4.173 to 0.821~kJ~mol$^{-1}$, a reduction of approximately 80\%. The relative-energy RMSE decreases from 4.857 to 1.036~kJ~mol$^{-1}$. Thus the basis improvement is not merely a common shift in absolute binding. A systematic relative bias remains: QZVPP underestimates eleven of the twelve energy differences, with a mean signed error of $-0.720$~kJ~mol$^{-1}$. For XIII, $\Delta E$ changes from $-2.994$ at TZVPP to $+0.692$~kJ~mol$^{-1}$ at QZVPP, compared with $+2.12$~kJ~mol$^{-1}$ from DMC. The QZVPP relative residual is consequently $-1.428$~kJ~mol$^{-1}$ despite the small absolute residual. The QZVPP result still places II and XI slightly below Ih by 0.190 and 0.033~kJ~mol$^{-1}$, respectively, whereas the DMC central values place them above Ih. The largest relative residual is $-2.390$~kJ~mol$^{-1}$ for VII. These residuals distinguish improved typical energy differences from an exact ordering of nearly degenerate phases. The comparison concerns electronic energies at fixed geometries, not finite-temperature phase stability, which also requires vibrational and pressure contributions.
